\subsection{Indiscrete Topology on Minkowski Space}\label{subsec:indiscrete-topology-on-minkowski-space}
Other topologies are possible on $\mathbb{R}^4$, and as Haag and Kastler don't explicitly specify a topology, we can explore other possibilities.

For example, one can equip $\mathbb{R}^4$ with the \href{https://ncatlab.org/nlab/show/discrete+and+indiscrete+topology}{indiscrete topology}. The indiscrete topology is the topology in which the only open sets are the empty set and $\mathbb{R}^4$ itself.

In this case, as we require our $\mathbf{B}$ to be an open, proper subset of $\mathbb{R}^4$, the only possible value that $\mathbf{B}$ can take on is the empty set. (Note, we require $\mathbf{B}$ to be a proper subset instead of just a subset to avoid the inclusion of global observables. See Section I Paragraph 5 of \href{https://doi.org/10.1063/1.1704187}{Haag Kastler}.) As $\mathbf{B}$ is intended to be a region in which one performs a ``measurement'', requiring $\mathbf{B}$ be the empty set doesn't capture that intention at all. The empty set is too ``small''.  So, we can assume that the indiscrete topology isn't implicitly being used.

But this leaves open the question: Exactly which topology \textit{is} being used? The Euclidean topology?
